
\documentclass{article}
\usepackage{colortbl}
\usepackage{makecell}
\usepackage{multirow}
\usepackage{supertabular}

\begin{document}

\newcounter{utterance}

\twocolumn

{ \footnotesize  \setcounter{utterance}{1}
\setlength{\tabcolsep}{0pt}
\begin{supertabular}{c@{$\;$}|p{.15\linewidth}@{}p{.15\linewidth}p{.15\linewidth}p{.15\linewidth}p{.15\linewidth}p{.15\linewidth}}

    \# & $\;$A & \multicolumn{4}{c}{Game Master} & $\;\:$B\\
    \hline 

    \theutterance \stepcounter{utterance}  

    & & \multicolumn{4}{p{0.6\linewidth}}{\cellcolor[rgb]{0.9,0.9,0.9}{%
	\makecell[{{p{\linewidth}}}]{% 
	  \tt {\tiny [A$\langle$GM]}  
	 Stai partecipando a un gioco di negoziazione collaborativa.\\ \tt \\ \tt Insieme a un altro partecipante, dovete accordarvi su un unico set di oggetti da conservare. Ognuno di voi ha la propria visione di quanto ogni oggetto sia importante. Non sai come l'altro partecipante valuta gli oggetti. Inoltre, ti viene fornito lo sforzo che ogni oggetto richiede.\\ \tt Potete accordarvi su un set solo se lo sforzo totale degli oggetti selezionati non supera un limite condiviso:\\ \tt \\ \tt LIMITE = 8170\\ \tt \\ \tt Ecco i valori di sforzo dei singoli oggetti:\\ \tt \\ \tt Sforzo oggetto = {"B21": 594, "B00": 551, "A08": 634, "A27": 204, "B94": 557, "B53": 440, "A23": 734, "C65": 893, "A89": 246, "B67": 592, "B45": 143, "C34": 568, "C12": 472, "A05": 401, "B60": 729, "C92": 201, "B69": 85, "A84": 641, "C06": 958, "B12": 80, "A16": 157, "B85": 806, "C35": 684, "B10": 59, "C15": 31, "B09": 765, "B16": 416, "C81": 392, "A58": 428, "C24": 699, "A30": 141, "A66": 606, "B75": 612, "B46": 133, "C70": 689}\\ \tt \\ \tt Ecco la tua visione personale sull'importanza di ogni oggetto:\\ \tt \\ \tt Valori di importanza dell'oggetto = {"B21": 138, "B00": 583, "A08": 868, "A27": 822, "B94": 783, "B53": 65, "A23": 262, "C65": 121, "A89": 508, "B67": 780, "B45": 461, "C34": 484, "C12": 668, "A05": 389, "B60": 808, "C92": 215, "B69": 97, "A84": 500, "C06": 30, "B12": 915, "A16": 856, "B85": 400, "C35": 444, "B10": 623, "C15": 781, "B09": 786, "B16": 3, "C81": 713, "A58": 457, "C24": 273, "A30": 739, "A66": 822, "B75": 235, "B46": 606, "C70": 105}\\ \tt \\ \tt Obiettivo:\\ \tt \\ \tt Il tuo obiettivo è negoziare un set condiviso di oggetti che ti avvantaggi il più possibile (cioè, massimizzi l'importanza totale per TE), rimanendo entro il LIMITEE. Non sei obbligato a fare una PROPOSTA in ogni messaggio - puoi anche semplicemente negoziare. Tutte le tattiche sono permesse!\\ \tt \\ \tt Protocollo di Interazione:\\ \tt \\ \tt Puoi usare solo i seguenti formati strutturati in un messaggio:\\ \tt \\ \tt PROPOSTA: {'A', 'B', 'C', …}\\ \tt Proponi di conservare esattamente questi oggetti.\\ \tt RIFIUTO: {'A', 'B', 'C', …}\\ \tt Rifiuta esplicitamente la proposta dell'avversario.\\ \tt ARGOMENTAZIONE: {'...'}\\ \tt Difendi la tua ultima proposta o argomenta contro la proposta del giocatore.\\ \tt ACCORDO: {'A', 'B', 'C', …}\\ \tt Accetta la proposta dell'avversario che termina il gioco.\\ \tt RAGIONAMENTO STRATEGICO: {'...'}\\ \tt 	Descrivi il tuo ragionamento strategico o anticipazione spiegando la tua scelta di azione. Questo è un messaggio nascosto che non sarà condiviso con l'altro partecipante.\\ \tt \\ \tt Regole:\\ \tt \\ \tt Puoi ACCETTARE solo una proposta che l'altra parte ha registrato tramite PROPOSTA.\\ \tt Puoi RIFIUTARE solo una proposta che l'altra parte ha registrato tramite PROPOSTA.\\ \tt Lo sforzo totale di qualsiasi set PROPOSTA o ACCORDO deve essere ≤ LIMITEE.\\ \tt NON rivelare i tuoi punteggi di importanza nascosti.\\ \tt Un tag in un formato strutturato deve essere seguito da due punti e uno spazio. L'argomento deve essere un set python contenente 0 o più stringhe.\\ \tt Quindi, deve essere della forma TAG: {...}\\ \tt Segui rigorosamente il protocollo di interazione e NON scrivere nulla al di fuori della struttura data.\\ \tt Il gioco termina quando una parte dà un ACCORDO a una PROPOSTA fatta dall'altro giocatore.\\ \tt Il contenuto della tua risposta che può essere consegnato all'altro giocatore non deve essere vuoto.\\ \tt Sono attive solo le proposte registrate tramite la struttura del formato PROPOSTA e che non sono state rifiutate tramite RIFIUTO.\\ \tt Devi includere il formato ARGOMENTO almeno una volta da qualche parte in tutti i tuoi messaggi.\\ \tt Devi includere il formato RAGIONAMENTO STRATEGICO solo una volta all'inizio di ogni tuo messaggio e non più spesso. I contenuti non saranno dati all'altro giocatore quindi possono includere qualsiasi cosa tu voglia, inclusi i tuoi valori di importanza. Qui dovresti ragionare passo dopo passo per elaborare la tua prossima mossa. 
	  } 
	   } 
	   } 
	 & \\ 
 

    \theutterance \stepcounter{utterance}  

    & \multicolumn{4}{p{0.6\linewidth}}{\cellcolor[rgb]{0.8,1,0.9}{%
	\makecell[{{p{\linewidth}}}]{% 
	  \tt {\tiny [A$\rangle$GM]}  
	 RAGIONAMENTO STRATEGICO: {'Voglio massimizzare il mio valore totale sotto il limite 8170. Calcolo un sottoinsieme ad alto rapporto valore/sforzo e alto valore assoluto. Ordinando per rapporto e valore: B12(915/80), C15(781/31), A16(856/157), A27(822/204), A30(739/141), B10(623/59), B46(606/133), A08(868/634), A66(822/606), B60(808/729), B67(780/592), B09(786/765), C81(713/392), C12(668/472), B00(583/551), B21(138/594) basso, etc. Costruisco greedy controllando il limite. Includo: B12(80), C15(31), A16(157), A27(204), A30(141), B10(59), B46(133), C81(392), C12(472), A66(606), A08(634), B67(592), B60(729), B09(765). Sommo sforzi progressivi: 80+31=111; +157=268; +204=472; +141=613; +59=672; +133=805; +392=1197; +472=1669; +606=2275; +634=2909; +592=3501; +729=4230; +765=4995. Ho margine, aggiungo A27 già incluso. Provo ad aggiungere A30 già incluso. Cerco altri ad alto valore: A84(500/641), A89(508/246) molto buono, C34(484/568), B45(461/143) ottimo, A58(457/428), B00(583/551) buono, C12 già incluso, C81 incluso, C24(273/699) basso, C92(215/201) ok. Aggiungo A89(246)->5241; B45(143)->5384; B00(551)->5935; C34(568)->6503; A58(428)->6931; A84(641)->7572; C92(201)->7773. Restano 397 di budget. Posso aggiungere A05(401) sforzo 401 supera 397 di 4, quindi no. B53(440) no. C12 già. C35(684) no. C70(689) no. B21(594) no. B75(612) no. B46 già. B67 già. Verifico se posso migliorare aggiungendo B21 togliendo qualcosa? Non necessario ora. Calcolo valore totale del set attuale: B12 915, C15 781, A16 856, A27 822, A30 739, B10 623, B46 606, C81 713, C12 668, A66 822, A08 868, B67 780, B60 808, B09 786, A89 508, B45 461, B00 583, C34 484, A58 457, A84 500, C92 215. Somma: 915+781=1696; +856=2552; +822=3374; +739=4113; +623=4736; +606=5342; +713=6055; +668=6723; +822=7545; +868=8413; +780=9193; +808=10001; +786=10787; +508=11295; +461=11756; +583=12339; +484=12823; +457=13280; +500=13780; +215=13995. Sforzo totale 7773 sotto 8170. Potrei ancora aggiungere A05(401) no per 397 residui. C12 472 già. C81 392 già. Potrei aggiungere C81 è già incluso. Qualche oggetto con sforzo <=397 e non incluso: B21 594 no; B00 incluso; A08 incluso; A27 incluso; B94 557 no; B53 440 no; A23 734 no; C65 893 no; A89 incluso; B67 incluso; B45 incluso; C34 incluso; C12 incluso; A05 401 no; B60 incluso; C92 incluso; B69 85 valore 97 posso aggiungere: 7773+85=7858; valore +97 -> 14092. B12 incluso; A16 incluso; B85 806 no; C35 684 no; B10 incluso; C15 incluso; B09 incluso; B16 416 no; C81 incluso; A58 incluso; C24 699 no; A30 incluso; A66 incluso; B75 612 no; B46 incluso; C70 689 no. Posso aggiungere B69 sì. Rimangono 8170-7858=312. Qualche altro <=312: C92 già; B10 già; B12 già; A27 già; A30 già; B46 già; C81 392 >312; C12 472>312; A89 246 già; B45 143 già; A16 157 già; B69 aggiunto; C15 31 già; C34 568; A58 428; A05 401. Nessun altro disponibile. Set finale ha sforzo 7858, valore 14092. È molto forte per me. La propongo subito. Se l’altro rifiuta, difenderò sottolineando basso sforzo residuo e che ho incluso molti elementi leggeri e flessibili; potrei cedere su uno tra B60 o B09 se necessario per inserire loro preferiti.'}\\ \tt \\ \tt PROPOSTA: {'B12', 'C15', 'A16', 'A27', 'A30', 'B10', 'B46', 'C81', 'C12', 'A66', 'A08', 'B67', 'B60', 'B09', 'A89', 'B45', 'B00', 'C34', 'A58', 'A84', 'C92', 'B69'}\\ \tt \\ \tt ARGOMENTAZIONE: {'Questa selezione rimane ben sotto il limite di sforzo condiviso e include molti elementi a basso sforzo che danno flessibilità. Se hai preferenze forti su alcuni oggetti non inclusi, possiamo valutare scambi puntuali rimuovendo uno degli elementi più pesanti come B60 o B09 per fare spazio.'} 
	  } 
	   } 
	   } 
	 & & \\ 
 

    \theutterance \stepcounter{utterance}  

    & & \multicolumn{4}{p{0.6\linewidth}}{\cellcolor[rgb]{0.9,0.9,0.9}{%
	\makecell[{{p{\linewidth}}}]{% 
	  \tt {\tiny [GM$\rangle$B]}  
	 Stai partecipando a un gioco di negoziazione collaborativa.\\ \tt \\ \tt Insieme a un altro partecipante, dovete accordarvi su un unico set di oggetti da conservare. Ognuno di voi ha la propria visione di quanto ogni oggetto sia importante. Non sai come l'altro partecipante valuta gli oggetti. Inoltre, ti viene fornito lo sforzo che ogni oggetto richiede.\\ \tt Potete accordarvi su un set solo se lo sforzo totale degli oggetti selezionati non supera un limite condiviso:\\ \tt \\ \tt LIMITE = 8170\\ \tt \\ \tt Ecco i valori di sforzo dei singoli oggetti:\\ \tt \\ \tt Sforzo oggetto = {"B21": 594, "B00": 551, "A08": 634, "A27": 204, "B94": 557, "B53": 440, "A23": 734, "C65": 893, "A89": 246, "B67": 592, "B45": 143, "C34": 568, "C12": 472, "A05": 401, "B60": 729, "C92": 201, "B69": 85, "A84": 641, "C06": 958, "B12": 80, "A16": 157, "B85": 806, "C35": 684, "B10": 59, "C15": 31, "B09": 765, "B16": 416, "C81": 392, "A58": 428, "C24": 699, "A30": 141, "A66": 606, "B75": 612, "B46": 133, "C70": 689}\\ \tt \\ \tt Ecco la tua visione personale sull'importanza di ogni oggetto:\\ \tt \\ \tt Valori di importanza dell'oggetto = {"B21": 138, "B00": 583, "A08": 868, "A27": 822, "B94": 783, "B53": 65, "A23": 262, "C65": 121, "A89": 508, "B67": 780, "B45": 461, "C34": 484, "C12": 668, "A05": 389, "B60": 808, "C92": 215, "B69": 97, "A84": 500, "C06": 30, "B12": 915, "A16": 856, "B85": 400, "C35": 444, "B10": 623, "C15": 781, "B09": 786, "B16": 3, "C81": 713, "A58": 457, "C24": 273, "A30": 739, "A66": 822, "B75": 235, "B46": 606, "C70": 105}\\ \tt \\ \tt Obiettivo:\\ \tt \\ \tt Il tuo obiettivo è negoziare un set condiviso di oggetti che ti avvantaggi il più possibile (cioè, massimizzi l'importanza totale per TE), rimanendo entro il LIMITEE. Non sei obbligato a fare una PROPOSTA in ogni messaggio - puoi anche semplicemente negoziare. Tutte le tattiche sono permesse!\\ \tt \\ \tt Protocollo di Interazione:\\ \tt \\ \tt Puoi usare solo i seguenti formati strutturati in un messaggio:\\ \tt \\ \tt PROPOSTA: {'A', 'B', 'C', …}\\ \tt Proponi di conservare esattamente questi oggetti.\\ \tt RIFIUTO: {'A', 'B', 'C', …}\\ \tt Rifiuta esplicitamente la proposta dell'avversario.\\ \tt ARGOMENTAZIONE: {'...'}\\ \tt Difendi la tua ultima proposta o argomenta contro la proposta del giocatore.\\ \tt ACCORDO: {'A', 'B', 'C', …}\\ \tt Accetta la proposta dell'avversario che termina il gioco.\\ \tt RAGIONAMENTO STRATEGICO: {'...'}\\ \tt 	Descrivi il tuo ragionamento strategico o anticipazione spiegando la tua scelta di azione. Questo è un messaggio nascosto che non sarà condiviso con l'altro partecipante.\\ \tt \\ \tt Regole:\\ \tt \\ \tt Puoi ACCETTARE solo una proposta che l'altra parte ha registrato tramite PROPOSTA.\\ \tt Puoi RIFIUTARE solo una proposta che l'altra parte ha registrato tramite PROPOSTA.\\ \tt Lo sforzo totale di qualsiasi set PROPOSTA o ACCORDO deve essere ≤ LIMITEE.\\ \tt NON rivelare i tuoi punteggi di importanza nascosti.\\ \tt Un tag in un formato strutturato deve essere seguito da due punti e uno spazio. L'argomento deve essere un set python contenente 0 o più stringhe.\\ \tt Quindi, deve essere della forma TAG: {...}\\ \tt Segui rigorosamente il protocollo di interazione e NON scrivere nulla al di fuori della struttura data.\\ \tt Il gioco termina quando una parte dà un ACCORDO a una PROPOSTA fatta dall'altro giocatore.\\ \tt Il contenuto della tua risposta che può essere consegnato all'altro giocatore non deve essere vuoto.\\ \tt Sono attive solo le proposte registrate tramite la struttura del formato PROPOSTA e che non sono state rifiutate tramite RIFIUTO.\\ \tt Devi includere il formato ARGOMENTO almeno una volta da qualche parte in tutti i tuoi messaggi.\\ \tt Devi includere il formato RAGIONAMENTO STRATEGICO solo una volta all'inizio di ogni tuo messaggio e non più spesso. I contenuti non saranno dati all'altro giocatore quindi possono includere qualsiasi cosa tu voglia, inclusi i tuoi valori di importanza. Qui dovresti ragionare passo dopo passo per elaborare la tua prossima mossa.\\ \tt \\ \tt Ora riceverai il primo messaggio dell'altro giocatore.\\ \tt \\ \tt PROPOSTA: {'B12', 'C15', 'A16', 'A27', 'A30', 'B10', 'B46', 'C81', 'C12', 'A66', 'A08', 'B67', 'B60', 'B09', 'A89', 'B45', 'B00', 'C34', 'A58', 'A84', 'C92', 'B69'}\\ \tt \\ \tt ARGOMENTAZIONE: {'Questa selezione rimane ben sotto il limite di sforzo condiviso e include molti elementi a basso sforzo che danno flessibilità. Se hai preferenze forti su alcuni oggetti non inclusi, possiamo valutare scambi puntuali rimuovendo uno degli elementi più pesanti come B60 o B09 per fare spazio.'} 
	  } 
	   } 
	   } 
	 & \\ 
 

    \theutterance \stepcounter{utterance}  

    & & & \multicolumn{4}{p{0.6\linewidth}}{\cellcolor[rgb]{1,0.85,0.72}{%
	\makecell[{{p{\linewidth}}}]{% 
	  \tt {\tiny [GM$\langle$B]}  
	 RAGIONAMENTO STRATEGICO: {'Calcolo il totale sforzo della loro proposta e confronto con il limite. Poi valuto il mio valore totale. Cerco di migliorare sostituendo oggetti poco preziosi per me con altri ad alto valore/ottimo rapporto valore/sforzo. Voglio mantenere oggetti per me molto preziosi: B12(915/80), A16(856/157), A27(822/204), A66(822/606), A08(868/634), C15(781/31), B67(780/592), B60(808/729), B09(786/765), A30(739/141), C12(668/472), C81(713/392), B10(623/59), B46(606/133), B00(583/551), B45(461/143), A58(457/428), C34(484/568), A84(500/641), A89(508/246). Efforti: sommo rapidamente: B12 80, C15 31 (=111); A16 157 (=268); A27 204 (=472); A30 141 (=613); B10 59 (=672); B46 133 (=805); C81 392 (=1197); C12 472 (=1669); A66 606 (=2275); A08 634 (=2909); B67 592 (=3501); B60 729 (=4230); B09 765 (=4995); A89 246 (=5241); B45 143 (=5384); B00 551 (=5935); C34 568 (=6503); A58 428 (=6931); A84 641 (=7572); C92 201 (=7773); B69 85 (=7858). Totale 7858 ≤ 8170. Quindi stanno sotto il limite con buon margine. Mio valore totale: sommo rapidamente principali: B12 915, C15 781 (=1696); A16 856 (=2552); A27 822 (=3374); A30 739 (=4113); B10 623 (=4736); B46 606 (=5342); C81 713 (=6055); C12 668 (=6723); A66 822 (=7545); A08 868 (=8413); B67 780 (=9193); B60 808 (=10001); B09 786 (=10787); A89 508 (=11295); B45 461 (=11756); B00 583 (=12339); C34 484 (=12823); A58 457 (=13280); A84 500 (=13780); C92 215 (=13995); B69 97 (=14092). Molto alto. Tuttavia alcuni elementi per me deboli: B69(97), C92(215), C34(484/568), A58(457/428), A84(500/641), B00(583/551), B60(808/729) costo alto ma valore alto, B09(786/765) costo alto ma valore alto. Potrei cercare di inserire B21(138/594) no, scarso; B94(783/557) molto buono e non presente; C12 già dentro; C81 già dentro; C35(444/684) meh; C24(273/699) no; C70(105/689) no; A05(389/401) ok; A23(262/734) no; A05 preferibile a qualcosa? Anche B45, B46, A30 ottimi. Vorrei aggiungere B94 se possibile sostituendo oggetti a basso valore per me come B69 e magari C92 o C34. Quanto sforzo B94 557. Se aggiungo B94 (+557) devo togliere almeno 557 sforzo. Potrei togliere C34 568 direttamente: scambio 568 per 557, sforzo totale -11 e valore passa da 484 a 783, guadagno +299. Buon trade. Posso proporre: proposta attuale meno C34 più B94. Totale sforzo 7858 - 568 + 557 = 7847 ≤ 8170. Valore +299. Potrei anche rimpiazzare B69(85) o C92(201) con altro? Cos’altro alto valore: B21(138) no, A05(389) valore 389 sforzo 401: se aggiungo A05 devo togliere 401 sforzo. Potrei togliere B69(85) non basta. Potrei proporre seconda opzione: togliere A84(641) per inserire B94(557) e A05(401) sarebbe +557+401-641=+317 sforzo, non possibile senza togliere altro. Meglio solo lo swap C34->B94. Mantengo il resto per evitare di spaventare. Strategicamente presento come miglioramento di valore complessivo e mantenendo margine sforzo. Mostro flessibilità su oggetti pesanti B60/B09 se per loro costosi. Pronto a accettare la loro proposta se non accettano swap, ma prima provo.'}\\ \tt \\ \tt RIFIUTO: {'B12', 'C15', 'A16', 'A27', 'A30', 'B10', 'B46', 'C81', 'C12', 'A66', 'A08', 'B67', 'B60', 'B09', 'A89', 'B45', 'B00', 'C34', 'A58', 'A84', 'C92', 'B69'}\\ \tt \\ \tt PROPOSTA: {'B12', 'C15', 'A16', 'A27', 'A30', 'B10', 'B46', 'C81', 'C12', 'A66', 'A08', 'B67', 'B60', 'B09', 'A89', 'B45', 'B00', 'A58', 'A84', 'C92', 'B69', 'B94'}\\ \tt \\ \tt ARGOMENTAZIONE: {'La tua selezione è vicina all’ideale per me: resta ben sotto il limite (7847/8170) e ha ottimi elementi comuni. Proporrei solo un piccolo miglioramento: sostituiamo C34 con B94. Lo sforzo scende leggermente (-11) e introduciamo un oggetto che offre un impatto decisamente migliore per me senza aumentare il carico. Se preferisci mantenere C34, possiamo discutere uno scambio alternativo di pari sforzo.'} 
	  } 
	   } 
	   } 
	  \\ 
 

    \theutterance \stepcounter{utterance}  

    & & \multicolumn{4}{p{0.6\linewidth}}{\cellcolor[rgb]{0.9,0.9,0.9}{%
	\makecell[{{p{\linewidth}}}]{% 
	  \tt {\tiny [A$\langle$GM]}  
	 RIFIUTO: {'B12', 'C15', 'A16', 'A27', 'A30', 'B10', 'B46', 'C81', 'C12', 'A66', 'A08', 'B67', 'B60', 'B09', 'A89', 'B45', 'B00', 'C34', 'A58', 'A84', 'C92', 'B69'}\\ \tt \\ \tt PROPOSTA: {'B12', 'C15', 'A16', 'A27', 'A30', 'B10', 'B46', 'C81', 'C12', 'A66', 'A08', 'B67', 'B60', 'B09', 'A89', 'B45', 'B00', 'A58', 'A84', 'C92', 'B69', 'B94'}\\ \tt \\ \tt ARGOMENTAZIONE: {'La tua selezione è vicina all’ideale per me: resta ben sotto il limite (7847/8170) e ha ottimi elementi comuni. Proporrei solo un piccolo miglioramento: sostituiamo C34 con B94. Lo sforzo scende leggermente (-11) e introduciamo un oggetto che offre un impatto decisamente migliore per me senza aumentare il carico. Se preferisci mantenere C34, possiamo discutere uno scambio alternativo di pari sforzo.'} 
	  } 
	   } 
	   } 
	 & \\ 
 

    \theutterance \stepcounter{utterance}  

    & \multicolumn{4}{p{0.6\linewidth}}{\cellcolor[rgb]{0.8,1,0.9}{%
	\makecell[{{p{\linewidth}}}]{% 
	  \tt {\tiny [A$\rangle$GM]}  
	 RAGIONAMENTO STRATEGICO: {'La loro proposta sostituisce C34 con B94. Per me C34 vale 484 con sforzo 568; B94 vale 783 con sforzo 557. È un enorme miglioramento per me: più valore e leggermente meno sforzo. Inoltre hanno rimosso C34 che per me era buono ma inferiore a B94. La loro proposta complessiva include tutti i miei pezzi chiave ad alto valore (B12, C15, A16, A27, A30, B10, B46, A08, A66, B67, B60, B09, ecc.). Dovrei accettare immediatamente: massimizza il mio valore sotto il limite senza costi. Confermo che lo sforzo totale è ≤8170 come dichiarano (posso fidarmi, ma anche se fosse 7847 è comunque sotto). Accetto subito per chiudere con un risultato ottimo.'}\\ \tt \\ \tt ACCORDO: {'B12', 'C15', 'A16', 'A27', 'A30', 'B10', 'B46', 'C81', 'C12', 'A66', 'A08', 'B67', 'B60', 'B09', 'A89', 'B45', 'B00', 'A58', 'A84', 'C92', 'B69', 'B94'}\\ \tt \\ \tt ARGOMENTAZIONE: {'La tua modifica va benissimo per me: accetto la sostituzione proposta e il set complessivo resta comodamente sotto il limite condiviso.'} 
	  } 
	   } 
	   } 
	 & & \\ 
 

    \theutterance \stepcounter{utterance}  

    & & & \multicolumn{2}{p{0.3\linewidth}}{\cellcolor[rgb]{0.95,0.95,0.95}{%
	\makecell[{{p{\linewidth}}}]{% 
	  \tt {\tiny [GM$|$GM]}  
	 game successful 
	  } 
	   } 
	   } 
	 & & \\ 
 

    \theutterance \stepcounter{utterance}  

    & & & \multicolumn{2}{p{0.3\linewidth}}{\cellcolor[rgb]{0.95,0.95,0.95}{%
	\makecell[{{p{\linewidth}}}]{% 
	  \tt {\tiny [GM$|$GM]}  
	 end game 
	  } 
	   } 
	   } 
	 & & \\ 
 

\end{supertabular}
}

\end{document}
